% CHAPTER 2: LITERATURE REVIEW
% ======================================================================
{\fontsize{18}{22}\selectfont\textbf{2 Literature Review}}\par

\vspace{3mm}

\subhead{2.1 LLMs in Enterprise Applications}

Transformer-based LLMs have changed AI research and its uses. Vaswani et al. (2017)[9] introduced the transformer architecture. Since then, this field gain a rapid progress, large models, like GPT-4, Claude, DeepSeek has come out. These models performs very well in understanding natural languages and inference tasks. Its core relies on self attention mechanism, which can capture semantic associations within the scope of long texts. Through this, coherent text will be generated.

In some specialized organizions, LLMs are used to automatically manage knowledge intensive works, including report writing, document abstract generating, programming code making, specialized field questions answering and complex work flow managing. But if the company use ready-made models, they may face lots of problems. Bender et al. (2021) [11] call these models "Random Parrot". They believe that models can only learn statistical laws from training data, with no really understanding ability. So they may output seemingly reasonable answers, which is wrong actually. In regulated industries like bank, this is worrying because it may led to legal penalties and economic losses.

Financial services industry has a strict attitude on using AI. Banks must meet the regulation commands. They need to standardize data confidentiality, transaction audit, and explanatory automate decision. The Basel Committee on Banking Supervision has issued guidelines on AI in the financial industry, requiring transparency, resilience testing, and manual supervision of model outputs.

\subhead{2.2 Multi-Agent Systems and Orchestration}

Multi agent systems (MAS) originated from research focus on distributed artificial intelligence. Early systems used rule-based agents, which had fixed interaction rules. As LLMs are becoming more and more popular, a new architecture design based on intelligent agents has emerged. Under this architecture, each agent relies on language models to complete reasoning, planning, and interactive communication between agents.

Wu et al. (2023) proposed AutoGen [4]. This platform allows different LLM-driven agents to communicate, and work together to solve complex questions. AutoGen has a dynamic conversation machanism, which can allocate each agent a role, so that subtasks can be transferred and peer review can be used to optimize outputs. But this frame is unpredictable and its output may not be reliable, which is unbearable in commercial scenarios. On the opposite, Hong et al. (2024)[5] proposed MetaGPT. This frame embeds Standard Operating Procedures (SOPs) into a multi-agent collaboration system, drawing inspiration from the working mode of factory assembly lines. Virtual roles such as product managers, system architects, development engineers, and testers are set up, and agents collaborate through the transfer of standardized intermediate results.

These frameworks show promising results in research settings. However, there's still a gap between flexible, role-based multi-agent systems and production systems in banking environments. For process standardization, systems must present deterministic behaviors and tight security controls.

\subhead{2.3 Code Generation and Text-to-SQL}

Converting natural language questions into executable SQL statements has been a major challenge in the field of natural language processing for over 20 years. The early solution adopted a rule-based approach, relying on manually constructed grammar rules and domain specific knowledge systems. This type of method is effective in narrow domain scenarios, but it is difficult to adapt to diverse database structures.

The emergence of the Spider dataset is an important development in this field. Yu et al. (2018) [1] constructed this dataset. Spider is a cross domain large-scale test set that includes 10181 query examples from 138 domains and 200 databases. Its core is used to test the model's ability to handle unfamiliar database architectures. This demand is highly compatible with the banking industry scenario: financial databases will be frequently updated due to new product launches, regulatory reporting requirements, and market data changes.

\subhead{2.4 RAG}

Retrieval-Augmented Generation (RAG) addresses an important weakness of large language models: their answers are generated from static parameters and may therefore be outdated, unverifiable, or hallucinated. Lewis et al. (2020) [8] introduced the core RAG idea of retrieving external evidence before generation. Instead of asking the model to answer from memory, a RAG system first searches a knowledge base for relevant passages and then gives those passages to the model as grounded context.

A typical RAG workflow has two connected phases. The offline indexing phase parses source documents, splits them into retrieval-sized chunks, converts each chunk into an embedding vector, and stores both vector data and traceable metadata. The online query phase embeds the user's question, retrieves semantically similar chunks from a vector database, applies ranking and filtering rules, and then asks the LLM to generate an answer from the retrieved evidence. This design separates language generation from organizational knowledge storage, allowing generated responses to use current enterprise documents without retraining the model.

Enterprise RAG differs from a simple demonstration system because retrieved text may be sensitive. In banking, policies, risk-control manuals, product rules, and compliance documents are often department-scoped and security-level controlled. Therefore, the retrieval stage must be coupled with access control: a high-similarity chunk cannot be used as evidence unless the requesting user is allowed to read the corresponding document. Citation output is also essential because employees must be able to verify which clause or document supports an answer.

RAG technology is highly valuable in the banking sector, as regulatory rules, internal policies, and product documentation are typically stored in vast text repositories. Furthermore, this content is frequently updated, making it impractical to integrate directly into model weights. Pipitone and Alami (2024) introduced LegalBench-RAG [2], a benchmark specifically designed to evaluate RAG systems in the legal and financial domains; this benchmark focuses on clause-level precision rather than document-level recall. This project follows the same principle: the RAG Agent is designed around chunk-level retrieval, permission-safe filtering, and citation-based responses rather than document-level keyword search.

\subhead{2.5 Tool-Augmented Model}

LLMs have inherent limitations that prevent them from directly interacting with external systems or data stores. Tool-augmented language models effectively solve this problem by equipping LLMs with the capability to call APIs, query databases, and perform computational tasks. Yao et al. (2023) [10] introduced the ReAct paradigm, whic combines "Reasoning" and "Acting", demonstrates that integrating reasoning steps with actions enables LLMs to use external tools for tasks involving real-time data, mathmatical calculations, and environmental interactions.

In enterprise applications, the scope of tool usage not only consists of standard web searches, but also specialized tools tailored to specific organizational needs, such as scheduling, resource booking, internal data querying, and process automation. The reliability of tool invocation is a critical performance metric, particularly regarding user intent understanding, precise parameter extraction, and error recovery.

\subhead{2.6 Data Security and Local Deployment}

Such tool invocation, however, is often realised through external cloud APIs, which raises additional security concerns. Security is core consideration of enterprises adopting AI systems. Currently, a common practice involves accessing LLMs via cloud APIs. That is, external endpoints accepts and processes users' queries and context. This method obeys the data governance policies enforced in regulated sectors. As Karras et al. (2025)[6] noted, sensitive information in prompt payloads may be logged by providers, repurposed for model retraining, or leaked through inversion techniques. 

In the banking sector, the Hong Kong Monetary Authority (HKMA) requires the implementation of appropriate safeguards to protect customer information, and also requires that relevant institutions retain control over data access and processing. Similar requirements are found in the EU’s General Data Protection Regulation (GDPR), the Personal Information Protection Law (PIPL) of mainland China, and banking security regulations in other jurisdictions. Together, these regulatory frameworks require that confidential financial data be processed in controlled environments that provide full auditability, which is a requirement that cannot be met when data is sent to or processed by third-party services.

On-premises deployment faces a significant computational cost challenge in LLM inference, which Kwon et al. (2023) [7] addressed with the PagedAttention mechanism and the vLLM framework. This solution, inspired by operating system virtual memory management, employs specific memory management strategies to effectively alleviate GPU memory bottlenecks.

Karras et al. (2024) [6] provided a comprehensive review of LLM applications in cybersecurity contexts in the big data era. Their analysis synthesized findings from over 100 research papers covering vulnerability detection, incident response, threat intelligence, and secure code generation. Basically, they found that while LLMs show promise in automating many cybersecurity tasks---such as identifying vulnerabilities in source code and generating security patches---current models still struggle with problems like contextual understanding and high false-positive rates in real-world enterprise settings. For this project, these findings reinforce how important it is to adopt a defense-in-depth security architecture instead of relying on LLM reasoning alone.
\newpage

% ======================================================================
